% =====================================================================
%  CLAD — single authoritative notation module.
%
%  One corpus, one symbol table. Every symbol used anywhere in the book is
%  defined here and nowhere else. If you need a symbol that is not here, add it
%  here first and check it against the collision list below before using it.
%
%  RESOLUTION PRINCIPLE
%  --------------------
%  Operators, predicates, and functions get readable multi-letter names.
%  Single letters are reserved for sets and their elements.
%
%  This is not a style preference. The governance corpus and the epistemic
%  corpus were written independently and collided on ten symbols; most of those
%  collisions were operator-versus-set, and the principle dissolves them without
%  anyone having to give up a letter they were using well.
%
%  COLLISIONS RESOLVED (see NOTATION-MAP.md for the full table)
%  ------------------------------------------------------------
%    O      deontic obligation | assertion operator
%           -> \Require and \Assert. O is retired entirely.
%    F      deontic prohibition | material facts f
%           -> \Forbid. f keeps material facts.
%    P, p   set of prompts | a proposition
%           -> prompts become x in X. p, q keep propositions.
%    C, c   operational constraints | resolved context
%           -> WITHDRAWN: "constraints become r in R" is unusable. R already
%              carries compliance risk R(i), the component requirement set
%              R_hard/R_soft in g = (S, C, E, A, R), and ROC's deterministic
%              rules r in R. Applying it would trade one collision for three.
%           -> ACTUAL DECISION: the two senses are chapter-disjoint and stay so.
%              c is an operational constraint in the governance chapters and
%              their appendix; c is the resolved context {j, t_v, f, u} in the
%              grounding chapters. Verified: neither sense appears in the
%              other's chapters. Capital C is unused outside the governance
%              material, so C for constraint sets carries no collision at all.
%              Anything introducing one sense into the other's chapters must
%              name the dimensions instead of using the symbol -- see the
%              applicability definition in Appendix A, which does exactly that.
%    A      agents / audit record | accuracy indicator
%           -> \Audit (set \AuditSet), agents \Agents. A_i keeps accuracy.
%    gamma  governability class | target grounding rate
%           -> \gov. gamma keeps the rate target.
%    S      control surface | the deployed system
%           -> surfaces become sigma_x in Sigma. S keeps the system.
%    V      set of audit records | token vocabulary
%           -> \AuditSet. V keeps the vocabulary.
%    D      governance domains | defeater search
%           -> \Defeaters. D keeps domains.
%    Phi    property vocabulary | component guarantee
%           -> \guarantee. Phi keeps the vocabulary. (An internal CLAD
%              collision, fixed here too.)
%
%  \newcommand throughout, never \providecommand: a name clash with a package
%  should break the build rather than silently resolve to someone else's macro.
% =====================================================================

% ---------------------------------------------------------------------
%  1. SETS AND ELEMENTS  (single letters, by the principle above)
% ---------------------------------------------------------------------
%   I           interactions
%   x in X      governed artifact — the assembled prompt or structured request
%   u           user input at runtime
%   M           the model: a parameterised generative function
%   S           the deployed system: M plus decoding, retrieval, policy, control
%   theta       inference configuration
%   o, o'       raw model output, delivered output
%   p, q        propositions / claims
%   c           resolved context, c = {j, t_v, f, u}
%   E           an evidentiary basis;  E_S  the system's evidential base
%   r in R      constraints (governance rules)
%   L           governance levels (enterprise > department > project)
%   D           governance domains (legal, security, data privacy, ...)
%   Phi         controlled vocabulary of atomic property identifiers
%   K           classifiers
%   T           time domain;  t_v valid time,  t_k transaction time
%   V           token vocabulary (V* = token sequences)
%   Sigma       control surfaces;  sigma_prompt, sigma_output, ...
%   V_proc      verification procedure
%   Inst        an institution
\newcommand{\Vocab}{\mathcal{V}}              % token vocabulary
\newcommand{\Levels}{L}
\newcommand{\Domains}{D}
\newcommand{\Props}{\Phi}                     % atomic property vocabulary
\newcommand{\Rules}{R}
\newcommand{\Surfaces}{\Sigma}
\newcommand{\Agents}{\mathcal{A}}             % was CLAD's A (agents)
\newcommand{\AuditSet}{\mathbb{A}}            % was CLAD's V (set of audit records)
\newcommand{\Inst}{I_{\mathrm{nst}}}          % institution; I alone is interactions

% ---------------------------------------------------------------------
%  2. DEONTIC OPERATORS
% ---------------------------------------------------------------------
%  These stay as O and F, reversing an earlier decision to rename them.
%
%  Two reasons. First, the collision that motivated the rename is gone: the
%  epistemic corpus's assertion operator is \Assert, so O is free. Second, and
%  decisively, the prompt-governance chapter engages the deontic literature
%  directly -- Ross's paradox, free-choice permission, why P_meta is not the
%  classical permission operator P, and why restricting to atomic O and F keeps
%  contradiction detection polynomial rather than NP-hard. That argument is not
%  decoration and it does not survive a rename: "P_meta is not the classical
%  deontic permission operator P" is incoherent if the obligation operator has
%  been renamed to Require.
%
%  The restricted two-operator fragment is a stated design decision with a
%  complexity justification, not an accident of notation.
%
%  Note on O: big-O complexity notation also appears (cost analysis in the
%  governed-retrieval chapter). The two are contextually unambiguous -- deontic O
%  takes a property identifier from Phi, big-O takes a complexity expression --
%  and both conventions coexist routinely.
\newcommand{\Obl}{\mathrm{O}}        % obligation: property must hold
\newcommand{\Prohib}{\mathrm{F}}     % prohibition: property must not hold
\newcommand{\Permitted}{\mathrm{P_{meta}}}  % governance annotation, NOT an operator
\newcommand{\Contradicts}{\operatorname{Contradicts}}
\newcommand{\Precedence}{\operatorname{precedence}}
\newcommand{\Tension}{\operatorname{Tension}}

% ---------------------------------------------------------------------
%  3. ASSERTION AND GROUNDING
% ---------------------------------------------------------------------
\newcommand{\Assert}{\operatorname{Assert}}    % was O_S: system asserts p
\newcommand{\Grd}{\operatorname{Grounded}}     % grounded assertion
\newcommand{\Grounded}{\operatorname{Grounded}_{E}} % grounded evidence set
\newcommand{\Sensitive}{\operatorname{Sensitive}}
% \gen, not \utter: the latter is speech-act notation and the prose never uses
% the word. Deliberately not \emit -- the Assert definition uses "emits" for the
% system-level event, and these two must not blur.
\newcommand{\gen}{\operatorname{gen}}

% --- the seven evidence conditions ---
\newcommand{\Authentic}{\operatorname{Authentic}}
\newcommand{\Versioned}{\operatorname{Versioned}}
\newcommand{\Authoritative}{\operatorname{Authoritative}}
\newcommand{\Applicable}{\operatorname{Applicable}}
\newcommand{\Supports}{\operatorname{Supports}}
\newcommand{\Sufficient}{\operatorname{Sufficient}}
\newcommand{\Traceable}{\operatorname{Traceable}}
\newcommand{\Entails}{\operatorname{Entails}}
\newcommand{\Current}{\operatorname{Current}}

% --- defeat ---
\newcommand{\NoDefeater}{\operatorname{NoDefeater}}
\newcommand{\DefeaterSearch}{\operatorname{DefeatSearch}} % was mathcal D; D is domains
\newcommand{\standing}{\operatorname{standing}}  % NOT alpha: alpha is the rate target
\newcommand{\GovQuery}{\mathcal{Q}}              % the governed retrieval

% ---------------------------------------------------------------------
%  4. TRUTH, CONTEXT, WARRANT
% ---------------------------------------------------------------------
\newcommand{\True}{\operatorname{True}}
\newcommand{\Resolved}{\operatorname{Resolved}}
\newcommand{\Evaluable}{\operatorname{Evaluable}}
\newcommand{\Verified}{\operatorname{Verified}}
\newcommand{\Authorized}{\operatorname{Authorized}}
\newcommand{\Warranted}{\operatorname{Warranted}}
\newcommand{\TrainedOn}{\operatorname{TrainedOn}}

% ---------------------------------------------------------------------
%  5. FAILURE CATEGORIES
% ---------------------------------------------------------------------
\newcommand{\Hall}{\operatorname{Hallucination}}
\newcommand{\Err}{\operatorname{Error}}
\newcommand{\Misapplied}{\operatorname{Misapplied}}
\newcommand{\Contested}{\operatorname{Contested}}
\newcommand{\Abstain}{\operatorname{Abstain}}

% ---------------------------------------------------------------------
%  6. GOVERNANCE STRUCTURE  (from the CLAD meta-framework)
% ---------------------------------------------------------------------
\newcommand{\gov}{\operatorname{gov}}          % was gamma(S): governability class
\newcommand{\guarantee}{\operatorname{guar}}   % was Phi(g): component guarantee
\newcommand{\Audit}{\operatorname{Audit}}      % was A_g(i,t)
\newcommand{\chain}{\operatorname{chain}}
\newcommand{\ver}{\operatorname{ver}}
\newcommand{\observe}{\operatorname{observe}}
\newcommand{\identity}{\operatorname{identity}}
\newcommand{\level}{\operatorname{level}}
\newcommand{\owner}{\operatorname{owner}}
\newcommand{\scope}{\operatorname{scope}}
\newcommand{\domainof}{\operatorname{domain}}
\newcommand{\decomp}{\operatorname{decomp}}
\newcommand{\risk}{\rho}                       % was R(i): compliance risk

% ---------------------------------------------------------------------
%  7. MODEL PERFORMANCE AND MEASUREMENT
% ---------------------------------------------------------------------
\newcommand{\Reliable}{\operatorname{Reliable}}
\newcommand{\LinguisticPlausibility}{\operatorname{LinguisticPlausibility}}
\newcommand{\SemanticCoherence}{\operatorname{SemanticCoherence}}
\newcommand{\BenchmarkAccuracy}{\operatorname{BenchmarkAccuracy}}
\newcommand{\Accuracy}{\operatorname{Accuracy}}
\newcommand{\GroundingRate}{\operatorname{GroundingRate}}

% ---------------------------------------------------------------------
%  8. RELATIONS AND SHORTHAND
% ---------------------------------------------------------------------
\newcommand{\nentails}{\;\not\Rightarrow\;}
\newcommand{\governs}{\succ}                   % enterprise > department > project
\newcommand{\satisfies}{\vDash}
